\chapter{Retail Flow and Wholesaling}\label{ch:m1:retail-flow-and-wholesaling}

You open a brokerage application, tap ``buy'', and ten shares appear in your
account a quarter of a second later at a price a fraction of a cent better
than the one on the screen. You paid no commission. Your order never reached
a stock exchange: your broker sent it to a trading firm, which sold you the
shares from its own inventory and paid your broker for the privilege. Dealers
trading as principal away from exchanges account for roughly four in ten
shares traded in the United States
(\cref{dat:m1:us-equity-market-structure:where}), and retail orders handled
this way are a large part of that. This chapter explains why a
trading firm pays to trade against you, how the quality of what you received
is measured, and what the arrangement does to everyone else.

\section{The path of a retail order}

\begin{definition}[Wholesaler and internalisation]\label{def:m1:retail-flow-and-wholesaling:wholesaler}
A \emph{wholesaler}\index{wholesaler} is a market-making firm that receives
the marketable orders of retail brokers' customers and executes most of them
as principal, against its own account, away from any exchange.
\emph{Internalisation}\index{internalisation} is the execution of a client
order against the executing firm's own inventory instead of against other
investors' orders on a venue.
\end{definition}

\begin{omfigure}
\begin{tikzpicture}[font=\small,
  box/.style={draw=omDef,rounded corners=2pt,minimum width=2.6cm,minimum height=0.95cm,align=center,fill=omDef!4},
  ext/.style={draw=omExo,rounded corners=2pt,minimum width=2.6cm,minimum height=0.95cm,align=center,fill=omExo!6},
  flow/.style={-{Stealth[length=2mm]},thick,omDef},
  pay/.style={-{Stealth[length=2mm]},dashed,thick,omThm}]
\node[ext] (c) at (0,0) {Retail\\customer};
\node[box] (b) at (4.3,0) {Retail broker\\(agent)};
\node[box] (w) at (8.6,0) {Wholesaler\\(principal)};
\node[ext] (x) at (12.9,1.3) {Exchanges};
\node[ext] (d) at (12.9,-1.3) {Dark pools};
\draw[flow] (c) -- node[above=10pt,font=\footnotesize]{order} (b);
\draw[flow] (b) -- node[above=10pt,font=\footnotesize]{routed} (w);
\draw[flow] (w) -- (x);
\draw[flow] (w) -- (d);
\draw[pay] (w.south west) to[bend left=28] node[below,font=\footnotesize]{payment for order flow} (b.south east);
\draw[pay] (b.south west) to[bend left=28] node[below,font=\footnotesize]{zero commission} (c.south east);
\end{tikzpicture}

\omcaption{A retail market order. The \omterm{def:m1:retail-flow-and-wholesaling:wholesaler}{wholesaler} fills it from inventory at
the national best price or slightly better, pays the broker, and manages the
resulting position on exchanges and in \omterm{def:m1:us-equity-market-structure:dark}{dark pools}. Only that residual
hedging reaches a public order book.}\label{fig:m1:retail-flow-and-wholesaling:path}
\end{omfigure}

The broker remains an agent: it owes its customer best execution, and the
choice of \omterm{def:m1:retail-flow-and-wholesaling:wholesaler}{wholesaler} is its routing decision. The \omterm{def:m1:retail-flow-and-wholesaling:wholesaler}{wholesaler} is a principal:
it owes the customer nothing beyond the price it promised the broker it
would deliver, and competes with other \omterm{def:m1:retail-flow-and-wholesaling:wholesaler}{wholesalers} on that promise.

\begin{definition}[Payment for order flow and price improvement]\label{def:m1:retail-flow-and-wholesaling:pfof}
\emph{Payment for order flow}\index{payment for order flow} (PFOF) is a
payment by an executing firm to a broker in return for the broker routing
its customers' orders to that firm. \emph{Price improvement}\index{price improvement}
is the amount by which an execution is better than the national best quote
on the customer's side at the time: for a purchase, the national best offer
less the price paid.
\end{definition}

\begin{dated}{2026-09}{The scale of the arrangement, and the reports about it}\label{dat:m1:retail-flow-and-wholesaling:scale}
By a trade publication's compilation of brokers' public routing reports, US
retail brokers received about \$950 million of \omterm{def:m1:retail-flow-and-wholesaling:pfof}{payment for order flow} in the
second quarter of 2025 alone, in stocks and options together. Amendments to
the SEC's execution-quality rule (Rule 605), adopted on 6 March 2024, extend
monthly reporting to large \omterm{def:m1:the-sell-side:sell-side}{broker-dealers} and add a plain summary report;
their compliance date was moved from 14 December 2025 to 1 August 2026, with
price-improvement statistics against the best displayed price following in
November 2026.
\end{dated}

\section{Measuring what the customer received}

\begin{definition}[Effective spread and realised spread]\label{def:m1:retail-flow-and-wholesaling:effective}
For an order of side $\epsilon = \pm1$ executed at price $p$ when the mid was
$m_t$, the \emph{effective half-spread}\index{effective spread} is
$\epsilon(p - m_t)$: what the order paid relative to the mid. The
\emph{realised half-spread}\index{realised spread} at horizon $\tau$ is
$\epsilon(p - m_{t+\tau})$: what the liquidity provider kept once the price
had moved. Both are usually quoted doubled, as spreads, and in cents or basis
points.
\end{definition}

\begin{proposition}[The spread identity]\label{prop:m1:retail-flow-and-wholesaling:identity}
For every trade and every horizon,
\[
\underbrace{\epsilon(p - m_t)}_{\text{effective}}
\;=\;
\underbrace{\epsilon(p - m_{t+\tau})}_{\text{realised}}
\;+\;
\underbrace{\epsilon(m_{t+\tau} - m_t)}_{\text{price impact}} .
\]
Averaged over an order flow, the price impact term is the \omterm{def:m1:what-a-trading-firm-does:adverse-selection}{adverse selection}
that flow inflicts (\cref{def:m1:what-a-trading-firm-does:adverse-selection}),
and the \omterm{def:m1:retail-flow-and-wholesaling:effective}{realised spread} is the liquidity provider's gross revenue per share.
\end{proposition}
\begin{proof}
Add and subtract $\epsilon\,m_{t+\tau}$.
\end{proof}

The identity turns the central idea of \cref{ch:m1:what-a-trading-firm-does}
into something measurable from a file of trades and quotes. It also shows
what an execution-quality statistic can and cannot say: the \omterm{def:m1:retail-flow-and-wholesaling:effective}{effective spread}
is what the \emph{customer} paid; the split between the two terms on the
right describes the \emph{flow}, not the quality of the execution.

\begin{definition}[Rule 605 and Rule 606 reports]\label{def:m1:retail-flow-and-wholesaling:reports}
A \emph{Rule 605 report}\index{Rule 605 report} is the monthly public
disclosure by which a US market centre states, by stock and order type and
size, its execution speed, effective and \omterm{def:m1:retail-flow-and-wholesaling:effective}{realised spreads} and \omterm{def:m1:retail-flow-and-wholesaling:pfof}{price
improvement}. A \emph{Rule 606 report}\index{Rule 606 report} is the quarterly
public disclosure by which a broker states to which venues it routed its
customers' orders and what payments it received from each.
\end{definition}

\begin{example}[Reading one fill]\label{ex:m1:retail-flow-and-wholesaling:onefill}
The national best quote is $20.00 \times 20.02$; a customer's market purchase
of 100 shares is filled at 20.018. \omterm{def:m1:retail-flow-and-wholesaling:pfof}{Price improvement}: 0.2 cent a share, 20
cents on the order. \omterm{def:m1:retail-flow-and-wholesaling:effective}{Effective half-spread}: $20.018 - 20.010 = 0.8$ cent.
Five minutes later the mid is 20.013: the \omterm{def:m1:retail-flow-and-wholesaling:effective}{realised half-spread} is $20.018 -
20.013 = 0.5$ cent and the price impact 0.3 cent. The \omterm{def:m1:retail-flow-and-wholesaling:wholesaler}{wholesaler}'s gross
revenue on this fill was 50 cents, out of which it pays the broker and its
own costs.
\end{example}

\section{Why uninformed flow is valuable}

\begin{proposition}[What a wholesaler can afford to pay]\label{prop:m1:retail-flow-and-wholesaling:maxpay}
Let the quoted half-spread be $h$. A flow whose orders are informed with
probability $p$, each informed order being followed by a move $J$, and which
is given a \omterm{def:m1:retail-flow-and-wholesaling:pfof}{price improvement} of a fraction $\theta$ of $h$, yields an
expected \omterm{def:m1:retail-flow-and-wholesaling:effective}{realised half-spread} of $h(1-\theta) - pJ$. A \omterm{def:m1:retail-flow-and-wholesaling:wholesaler}{wholesaler} with
operating costs $\kappa$ per share breaks even paying the broker at most
\[
\pi_{\max} \;=\; h(1-\theta) - pJ - \kappa .
\]
\end{proposition}
\begin{proof}
\omterm{def:m1:retail-flow-and-wholesaling:effective}{Effective half-spread} $h(1-\theta)$ less expected price impact $pJ$ is the
\omterm{def:m1:retail-flow-and-wholesaling:effective}{realised half-spread} by \cref{prop:m1:retail-flow-and-wholesaling:identity};
subtract costs.
\end{proof}

\begin{example}[The same cent, two flows]\label{ex:m1:retail-flow-and-wholesaling:twoflows}
Take $h = 1$ cent, $J = 3$ cents, $\kappa = 0.10$ cent. On an exchange,
where 35\% of the orders hitting a quote are informed, a \omterm{def:m1:what-a-trading-firm-does:market-maker}{market maker} earns
$1 - 1.05 < 0$ before costs: the quoted spread is, in equilibrium, barely
enough. For retail flow with $p = 5\%$ and $\theta = 20\%$: $\pi_{\max} =
0.80 - 0.15 - 0.10 = 0.55$ cent. The \omterm{def:m1:retail-flow-and-wholesaling:wholesaler}{wholesaler} can pay the broker 0.25 cent
and keep 0.30 (\cref{fig:m1:retail-flow-and-wholesaling:halfspread}). Nothing
here depends on speed or on being cleverer than the customer: the value lies
entirely in the \emph{segmentation} of orders by how much they know.
\end{example}

\begin{omfigure}
\begin{tikzpicture}
\begin{axis}[width=9cm,height=4.6cm,xbar,bar width=8pt,
  xlabel={cents per share, out of a quoted half-spread of 1 cent},
  ytick=data,yticklabels from table={figdata/markets-1/10-retail-flow-and-wholesaling/halfspread.csv}{label},
  yticklabel style={font=\small},tick label style={font=\small},label style={font=\small},
  y dir=reverse,xmin=0,xmax=0.4,xtick={0,0.1,0.2,0.3,0.4},enlarge y limits=0.15,grid=major,grid style={gray!25},
  xticklabel style={/pgf/number format/fixed,/pgf/number format/precision=1},
  nodes near coords,nodes near coords style={font=\footnotesize,/pgf/number format/fixed,/pgf/number format/precision=2},
  table/col sep=comma]
\addplot[fill=omDef!60,draw=omDef] table[x=cents,y=k]{figdata/markets-1/10-retail-flow-and-wholesaling/halfspread.csv};
\end{axis}
\end{tikzpicture}

\omcaption{Where one cent of quoted half-spread goes when the order is a
retail one. Parameters are those of
\cref{ex:m1:retail-flow-and-wholesaling:twoflows} and are illustrative. Data:
the chapter's script.}\label{fig:m1:retail-flow-and-wholesaling:halfspread}
\end{omfigure}

\begin{omfigure}
\begin{tikzpicture}
\begin{axis}[width=10.5cm,height=4.8cm,xmode=log,
  xlabel={time after the trade (seconds)},ylabel={realised half-spread (cents)},
  tick label style={font=\small},label style={font=\small},
  xmin=0.1,xmax=300,ymin=-0.2,ymax=1.1,grid=major,grid style={gray!25},
  legend columns=2,legend style={font=\footnotesize,at={(0.5,-0.36)},anchor=north,draw=none,fill=none,/tikz/every even column/.append style={column sep=8pt}},
  table/col sep=comma]
\addplot[omDef,very thick,mark=*,mark size=1.4pt] table[x=horizon_s,y=retail_realised]{figdata/markets-1/10-retail-flow-and-wholesaling/markout.csv};
\addplot[omThm,very thick,mark=*,mark size=1.4pt] table[x=horizon_s,y=exchange_realised]{figdata/markets-1/10-retail-flow-and-wholesaling/markout.csv};
\legend{retail flow (with 20\% improvement),flow hitting exchange quotes}
\end{axis}
\end{tikzpicture}

\omcaption{What the liquidity provider keeps, against how long after the
trade it is measured (logarithmic horizontal axis). Both flows start near the
quoted half-spread; five minutes later the exchange flow has taken all of it
back and the retail flow has left two thirds. Data: the tutorial's
simulation, 400\,000 trades per curve.}\label{fig:m1:retail-flow-and-wholesaling:markout}
\end{omfigure}

\begin{omfigure}
\begin{tikzpicture}
\begin{axis}[width=10.5cm,height=4.4cm,
  xlabel={share of informed orders in the flow (\%)},ylabel={$\pi_{\max}$ (cents a share)},
  tick label style={font=\small},label style={font=\small},
  xmin=0,xmax=30,ymin=-0.3,ymax=0.8,grid=major,grid style={gray!25},table/col sep=comma]
\addplot[omDef,very thick] table[x=informed_pct,y=max_payment]{figdata/markets-1/10-retail-flow-and-wholesaling/maxpay.csv};
\addplot[black,dashed] coordinates {(0,0) (30,0)};
\addplot[only marks,mark=*,omThm] coordinates {(5,0.55)};
\node[omThm,font=\footnotesize,anchor=south west,fill=white,inner sep=1.5pt] at (axis cs:5.6,0.57) {retail flow of the example};
\end{axis}
\end{tikzpicture}

\omcaption{The most a \omterm{def:m1:retail-flow-and-wholesaling:wholesaler}{wholesaler} can pay for a flow, against how informed the
flow is, for a quoted half-spread of one cent, 20\% improvement, a 3-cent
informed move and costs of 0.10 cent. Beyond 23\% of informed orders the flow
has a negative price. Data: \cref{prop:m1:retail-flow-and-wholesaling:maxpay},
computed by the chapter's script.}
\end{omfigure}

\begin{remark}[Who pays]\label{rem:m1:retail-flow-and-wholesaling:whopays}
The retail customer is better off than on an exchange by the \omterm{def:m1:retail-flow-and-wholesaling:pfof}{price
improvement} and the absent commission. The cost falls elsewhere. Exchange
quotes are set for the flow that reaches exchanges; if the mildest orders are
removed from it, $p$ rises there and so does the spread
(\cref{rem:m1:us-equity-market-structure:free}). \omterm{def:m1:retail-flow-and-wholesaling:pfof}{Price improvement} is then
measured against a benchmark that the arrangement itself has widened. How
large this effect is, and whether order-by-order competition would serve
customers better than wholesaling, has long been argued before the SEC
without a settled answer.
\end{remark}

\section{When the arrangement is tested}

\begin{example}[An enforcement case]\label{ex:m1:retail-flow-and-wholesaling:case}
In December 2020 the SEC charged Robinhood Financial with misleading
customers about its revenue sources and with failing its duty of best
execution; the firm agreed to pay \$65 million without admitting or denying
the findings. The order found that between 2015 and late 2018 the broker's
customer communications omitted \omterm{def:m1:retail-flow-and-wholesaling:pfof}{payment for order flow}, its largest source
of revenue, and that it had accepted unusually high payments in exchange for
lower \omterm{def:m1:retail-flow-and-wholesaling:pfof}{price improvement}: its customers' executions were \$34.1 million worse
than at competing brokers, \emph{even after} allowing for the commissions
they did not pay.
\end{example}

The case states the conflict precisely: \cref{prop:m1:retail-flow-and-wholesaling:maxpay}
fixes the total, $h - pJ - \kappa$, and leaves its division between the
customer ($\theta h$), the broker ($\pi$) and the \omterm{def:m1:retail-flow-and-wholesaling:wholesaler}{wholesaler} to negotiation
between the last two. The customer is not at the table; the Rule 605 and 606
reports, and the broker's best-execution duty, are there on its behalf.

The week of 28 January 2021 tested the chain at a different link: not
pricing but \emph{plumbing}. The purchase restrictions that made headlines
came from clearing-house collateral (\cref{ch:m1:clearing-and-settlement}),
not from the \omterm{def:m1:retail-flow-and-wholesaling:wholesaler}{wholesalers}; the SEC staff's report on the episode concluded
that ``it was the positive sentiment, not the buying-to-cover, that sustained
the weeks-long price appreciation'' of the stock at its centre.
It is the first case study of \cref{ch:m1:stress-case-studies}.

\section{Tutorial: measuring execution quality}\label{tut:m1:retail-flow-and-wholesaling:execq}

\begin{tutorial}
\textbf{Goal.} Compute \omterm{def:m1:retail-flow-and-wholesaling:effective}{effective spread}, \omterm{def:m1:retail-flow-and-wholesaling:effective}{realised spread} and price impact on
two simulated order flows and reproduce
\cref{fig:m1:retail-flow-and-wholesaling:markout}. \textbf{End state:} a
\omterm{def:m1:retail-flow-and-wholesaling:effective}{realised half-spread} after five minutes of 0.65 cent for retail flow and
about zero for exchange flow.
\begin{tutsteps}
\item \textbf{The flows.} Orders are informed with probability 5\% (retail)
or 35\% (exchange); an informed order is followed by a 3-cent move that
unfolds over about a minute; noise is added at every horizon.
\omcode{code/markets-1/10-retail-flow-and-wholesaling/python/execq.py}{18}{38}{Two flows that differ only in how much they know and in the price improvement they are given.}\label{lst:m1:retail-flow-and-wholesaling:flows}
\item \textbf{The statistics}: three means, and the identity between them.
\omcode{code/markets-1/10-retail-flow-and-wholesaling/python/execq.py}{41}{46}{Effective, realised, impact.}\label{lst:m1:retail-flow-and-wholesaling:stats}
\item \textbf{Run.} \omterm{def:m1:retail-flow-and-wholesaling:effective}{Effective half-spreads} are 0.80 and 1.00 cent. At 300
seconds the price impacts are 0.15 and 1.05 cents, the products $pJ$; the
\omterm{def:m1:retail-flow-and-wholesaling:effective}{realised half-spreads} are what remains.
\item \textbf{Choose the horizon.} At 100 milliseconds both flows look
equally profitable; the difference appears only over tens of seconds. A
horizon shorter than the time over which information reaches the price
flatters every flow.
\end{tutsteps}
\textbf{What to change next.} Give retail orders in one stock $p = 30\%$ (a
``meme'' stock in a frenzy) and recompute $\pi_{\max}$. Then make the noise
ten times larger and see how many trades the means need before the two flows
can be told apart.
\end{tutorial}

\section{Build: the execution-quality report}\label{bld:m1:retail-flow-and-wholesaling:report}

\begin{build}
\bfield{Purpose} Every execution of the miniature firm, whether it provides
or takes liquidity, is scored the same way. The report is the firm's own
Rule 605 table and the raw material of the transaction-cost analysis of One
Quant Book 10.
\bfield{Interface} \texttt{Fill(ts, symbol, side, quantity, price, bid, ask)}
with the quote at execution; a \texttt{mid\_at(symbol, ts)} callback;
\texttt{report(fills, horizons, by=(symbol, size\_bucket))} returning,
per group: number of fills and shares, share-weighted effective, realised and
impact half-spreads in basis points, \omterm{def:m1:retail-flow-and-wholesaling:pfof}{price improvement} in cents, and the
fraction of shares executed at, inside and outside the quote.
\bfield{Rules} Integer prices; the identity of
\cref{prop:m1:retail-flow-and-wholesaling:identity} must hold exactly per
fill. A fill with a crossed or missing quote is excluded and counted
separately, never silently dropped. Size buckets follow the customary order
size ranges (1--99, 100--499, 500--1\,999, 2\,000--4\,999, 5\,000 and more).
\bfield{Acceptance tests} \texttt{code/firm/execquality/tests/}: the
single-fill example of this chapter; the identity on random fills; weighting
by shares, not by fills; exclusions counted.
\bfield{Stretch} Add the markout curve: \omterm{def:m1:retail-flow-and-wholesaling:effective}{realised spread} at a list of
horizons, with a standard error per point.
\end{build}

\begin{omsources}
\item US Securities and Exchange Commission, \emph{SEC Charges Robinhood
Financial With Misleading Customers About Revenue Sources and Failing to
Satisfy Duty of Best Execution}, press release 2020-321, 17 December 2020,
and administrative order 33-10906.
\item US Securities and Exchange Commission, \emph{Disclosure of Order
Execution Information}, adopting release of 6 March 2024; \emph{Extension of
Compliance Date}, \emph{Federal Register}, 2 October 2025.
\item US Securities and Exchange Commission, \emph{Staff Report on Equity and
Options Market Structure Conditions in Early 2021}, October 2021.
\item Global Trading, ``Robinhood, Schwab led retail order flow payment
bonanza'', 2025 (compilation of Rule 606 reports).
\item R.~Huang and H.~Stoll, ``Dealer versus auction markets: a paired
comparison of execution costs on NASDAQ and the NYSE'', \emph{Journal of
Financial Economics} 41 (1996) --- effective and realised spreads.
\item R.~Battalio, S.~Corwin and R.~Jennings, ``Can brokers have it all? On
the relation between make-take fees and limit order execution quality'',
\emph{Journal of Finance} 71 (2016).
\end{omsources}

\section{Exercises}

\begin{exercise}[$\star$]\label{exo:m1:retail-flow-and-wholesaling:1}
The quote is $35.40 \times 35.44$. A customer sells 200 shares at 35.405.
Compute the \omterm{def:m1:retail-flow-and-wholesaling:pfof}{price improvement} per share and in dollars, and the \omterm{def:m1:retail-flow-and-wholesaling:effective}{effective
half-spread}.
\end{exercise}

\begin{exercise}[$\star$]\label{exo:m1:retail-flow-and-wholesaling:2}
For the sale of the previous exercise the mid one minute later is 35.412.
Compute the \omterm{def:m1:retail-flow-and-wholesaling:effective}{realised half-spread} and the price impact, and check the
identity.
\end{exercise}

\begin{exercise}[$\star$]\label{exo:m1:retail-flow-and-wholesaling:3}
A broker receives 0.18 cent a share of \omterm{def:m1:retail-flow-and-wholesaling:pfof}{payment for order flow} on 60 million
shares a day. What is that per year (252 days)? Express it per customer if
the broker has 12 million funded accounts.
\end{exercise}

\begin{exercise}[$\star\star$]\label{exo:m1:retail-flow-and-wholesaling:4}
With $h = 1.5$ cents, $J = 4$ cents, $\kappa = 0.12$ cent and a \omterm{def:m1:retail-flow-and-wholesaling:pfof}{price
improvement} of 15\% of $h$: compute $\pi_{\max}$ for flows with $p = 4\%$ and
$p = 20\%$. At what $p$ does the \omterm{def:m1:retail-flow-and-wholesaling:wholesaler}{wholesaler} stop bidding for the flow?
\end{exercise}

\begin{exercise}[$\star\star$]\label{exo:m1:retail-flow-and-wholesaling:5}
Broker A takes 0.30 cent of payment and its customers get 10\% of the
half-spread as improvement; broker B takes 0.10 cent and its customers get
30\%. With $h = 1$ cent, how much better off per share is B's customer? If B
charges a commission of \$1 per order, for what order sizes is A cheaper?
\end{exercise}

\begin{exercise}[$\star\star$]\label{exo:m1:retail-flow-and-wholesaling:6}
A \omterm{def:m1:retail-flow-and-wholesaling:wholesaler}{wholesaler} measures \omterm{def:m1:retail-flow-and-wholesaling:effective}{realised spreads} of 0.70 cent at one second and 0.20
cent at five minutes on a broker's flow, with an \omterm{def:m1:retail-flow-and-wholesaling:effective}{effective half-spread} of
0.75. Compute the price impact at both horizons. What does the pattern say
about this broker's customers, and what will the \omterm{def:m1:retail-flow-and-wholesaling:wholesaler}{wholesaler} do?
\end{exercise}

\begin{exercise}[$\star\star\star$]\label{exo:m1:retail-flow-and-wholesaling:7}
\emph{Coding.} With \texttt{simulate} and \texttt{stats}, estimate for the
retail flow the standard error of the mean \omterm{def:m1:retail-flow-and-wholesaling:effective}{realised half-spread} at 300
seconds from $n = 10\,000$ trades, and deduce how many trades are needed to
distinguish $p = 5\%$ from $p = 8\%$ at two standard errors.
\end{exercise}

\begin{exercise}[$\star\star\star$]\label{exo:m1:retail-flow-and-wholesaling:8}
\emph{Find the flaw.} A \omterm{def:m1:retail-flow-and-wholesaling:wholesaler}{wholesaler}'s marketing states: ``We delivered \$1.2
billion of \omterm{def:m1:retail-flow-and-wholesaling:pfof}{price improvement} to retail investors last year.'' Give three
reasons why this number, even if exactly computed, does not measure the
benefit to those investors.
\end{exercise}

\section{Problem: Pricing Retail Flow}

\begin{problem}[{Weekend problem --- bidding for a broker's order flow}]\label{pb:m1:retail-flow-and-wholesaling:1}
You run the wholesaling desk of a market-making firm. A retail broker routing
40 million shares a day invites you to bid. Its flow is in stocks with an
average quoted half-spread of 1.2 cents. From a sample you estimate that 6\%
of its orders are informed, with an average subsequent move of 3.5 cents.
Your costs are 0.11 cent a share.

\medskip\noindent\textbf{Part I --- The economics of the flow.}
\begin{enumerate}
\item Compute the expected price impact per share.
\item With a \omterm{def:m1:retail-flow-and-wholesaling:pfof}{price improvement} of 20\% of the half-spread, give the
effective and the \omterm{def:m1:retail-flow-and-wholesaling:effective}{realised half-spread}.
\item Compute $\pi_{\max}$.
\item Give your daily gross revenue and daily margin if you pay 0.30 cent.
\item And per year of 252 days.
\end{enumerate}

\medskip\noindent\textbf{Part II --- The auction.}
The broker ranks \omterm{def:m1:retail-flow-and-wholesaling:wholesaler}{wholesalers} by \omterm{def:m1:retail-flow-and-wholesaling:pfof}{price improvement} delivered, and sets the
payment itself at 0.30 cent for all.
\begin{enumerate}[resume]
\item What is the largest improvement fraction $\theta$ at which you break
even?
\item A competitor has costs of 0.06 cent. What $\theta$ can it offer?
\item You can hedge residual inventory more cheaply than it can, which lowers
your effective price impact by 0.08 cent. Who wins now?
\item Why does the broker route to several \omterm{def:m1:retail-flow-and-wholesaling:wholesaler}{wholesalers} at once and reshuffle
shares monthly?
\end{enumerate}

\medskip\noindent\textbf{Part III --- The flow changes.}
\begin{enumerate}[resume]
\item During a speculative frenzy the informed share of orders in twenty
stocks rises to 25\% and the move to 6 cents. Compute the \omterm{def:m1:retail-flow-and-wholesaling:effective}{realised
half-spread} in those stocks at $\theta = 20\%$ and unchanged spreads.
\item Those stocks are 15\% of the broker's volume. Give your new daily
margin, paying 0.30 cent.
\item Spreads in those stocks widen to a half-spread of 3 cents. Redo
question 10.
\item List three things you can do, within your agreement with the broker,
when one segment of a flow turns toxic.
\end{enumerate}

\medskip\noindent\textbf{Part IV --- The other parties.}
\begin{enumerate}[resume]
\item What does the broker earn per year from your payment?
\item What do its customers receive per year in \omterm{def:m1:retail-flow-and-wholesaling:pfof}{price improvement}?
\item The broker could instead send all orders to an exchange charging a
taker fee of 0.30 cent. What would the customers pay relative to the mid, and
what would the broker pay?
\item Compare the customers' all-in cost in the two arrangements. What has
been assumed about the exchange's quoted spread?
\item Under \cref{rem:m1:retail-flow-and-wholesaling:whopays}, how might that
assumption fail if all retail flow moved to exchanges?
\item State the \emph{named result}: the maximum payment per share at which
you break even on this flow, in cents.
\item In one sentence: what exactly is the \omterm{def:m1:retail-flow-and-wholesaling:wholesaler}{wholesaler} buying?
\end{enumerate}
\end{problem}

\section{Interview questions}

\begin{interviewq}[$\star$ \iqroles{trader, researcher, developer}]\label{iq:m1:retail-flow-and-wholesaling:1}
What is \omterm{def:m1:retail-flow-and-wholesaling:pfof}{payment for order flow}, and why would a \omterm{def:m1:what-a-trading-firm-does:market-maker}{market maker} pay for orders?
\end{interviewq}

\begin{interviewq}[$\star$ \iqroles{trader, researcher}]\label{iq:m1:retail-flow-and-wholesaling:2}
Define the \omterm{def:m1:retail-flow-and-wholesaling:effective}{effective spread} and the \omterm{def:m1:retail-flow-and-wholesaling:effective}{realised spread}. Which one measures the
customer's cost and which one the \omterm{def:m1:what-a-trading-firm-does:market-maker}{market maker}'s revenue?
\end{interviewq}

\begin{interviewq}[$\star\star$ \iqroles{researcher, mle}]\label{iq:m1:retail-flow-and-wholesaling:3}
You are given a year of a broker's fills with quotes. How do you decide
whether its flow is ``toxic'', and at which horizon do you measure?
\end{interviewq}

\begin{interviewq}[$\star\star$ \iqroles{trader, researcher}]\label{iq:m1:retail-flow-and-wholesaling:4}
A retail order and an institutional algorithm's child order are both for 100
shares at market. Why does a \omterm{def:m1:what-a-trading-firm-does:market-maker}{market maker} treat them differently, and how
can it tell them apart on an anonymous exchange?
\end{interviewq}

\begin{interviewq}[$\star\star$ \iqroles{researcher, bank}]\label{iq:m1:retail-flow-and-wholesaling:5}
Make the case that \omterm{def:m1:retail-flow-and-wholesaling:pfof}{payment for order flow} is good for retail investors, then
the case that it is bad for markets.
\end{interviewq}

\begin{interviewq}[$\star\star\star$ \iqroles{researcher, trader}]\label{iq:m1:retail-flow-and-wholesaling:6}
Your wholesaling desk's \omterm{def:m1:retail-flow-and-wholesaling:effective}{realised spread} on one broker's flow has halved over
six months while its \omterm{def:m1:retail-flow-and-wholesaling:effective}{effective spread} is unchanged. Give three hypotheses
and the data that would discriminate between them.
\end{interviewq}
